% Part:sets-functions-relations
% Chapter: size-of-sets
% Section: zig-zag

\documentclass[../../../include/open-logic-section]{subfiles}

\begin{document}

\olfileid{sfr}{siz}{zigzag}

\olsection{కాంటర్ జిగ్‌జాగ్ పద్ధతి}

\begin{explain}
ఇప్పటికే కొన్ని ``సులభమైన'' లెక్కింపులను పరిశీలించాం. ఇప్పుడు కాస్త
కష్టమైనదాన్ని చూద్దాం. సహజ సంఖ్యల జతల సమితిని పరిశీలించండి;
\oliflabeldef{sfr:set:pai:sec}{దీనిని \olref[set][pai]{sec}లో ఇలా
నిర్వచించాం:}{దీని నిర్వచనం:}
\[
\Nat \times \Nat = \Setabs{\tuple{n,m}}{n,m \in \Nat}
\]
ఈ క్రమయుగ్మాలను ఒక \emph{పట్టిక}లో ఇలా అమర్చవచ్చు:
\[
\begin{array}{ c | c | c | c | c | c}
& \mathbf 0 & \mathbf 1 & \mathbf 2 & \mathbf 3 & \dots \\
\hline
\mathbf 0 & \tuple{0,0} & \tuple{0,1} & \tuple{0,2} & \tuple{0,3} & \dots \\
\hline
\mathbf 1 & \tuple{1,0} & \tuple{1,1} & \tuple{1,2} & \tuple{1,3} & \dots \\
\hline
\mathbf 2 & \tuple{2,0} & \tuple{2,1} & \tuple{2,2} & \tuple{2,3} & \dots \\
\hline
\mathbf 3 & \tuple{3,0} & \tuple{3,1} & \tuple{3,2} & \tuple{3,3} & \dots \\
\hline
\vdots & \vdots & \vdots & \vdots & \vdots & \ddots\\
\end{array}
\]
$\Nat \times \Nat$లోని ప్రతి క్రమయుగ్మం పట్టికలో కచ్చితంగా ఒక్కసారి
కనిపిస్తుంది. ప్రత్యేకించి, $\tuple{n,m}$ అనేది $n$వ అడ్డువరుసలో,
$m$వ నిలువువరుసలో కనిపిస్తుంది. కానీ ఇలాంటి పట్టికలోని మూలకాలను
``ఏకమితీయ'' జాబితాగా ఎలా అమర్చాలి? కింది పట్టికలోని నమూనా ఒక పద్ధతిని
చూపుతుంది (వాస్తవానికి మరెన్నో మార్గాలు ఉన్నాయి):
\[
\begin{array}{ c | c | c | c | c | c | c}
& \mathbf 0 & \mathbf 1 & \mathbf 2 & \mathbf 3 & \mathbf 4 &\dots \\
\hline
\mathbf 0 & 0  & 1& 3 & 6& 10 &\ldots \\
\hline
\mathbf 1 &2 & 4& 7 & 11 & \dots &\ldots \\
\hline
\mathbf 2 & 5 & 8 & 12 & \ldots & \dots&\ldots \\
\hline
\mathbf 3 & 9 & 13 & \ldots & \ldots & \dots & \ldots \\
\hline
\mathbf 4 & 14 & \ldots & \ldots & \ldots & \dots & \ldots \\
\hline
\vdots & \vdots & \vdots & \vdots & \vdots&\ldots & \ddots\\
\end{array}
\]\noindent
ఈ నమూనాను \emph{కాంటర్ జిగ్‌జాగ్ పద్ధతి} అంటారు. ఇది
$\Nat \times \Nat$ను ఇలా లెక్కిస్తుంది:
\[
\tuple{0,0}, \tuple{0,1}, \tuple{1,0}, \tuple{0,2}, \tuple{1,1},
\tuple{2,0}, \tuple{0,3}, \tuple{1,2}, \tuple{2,1}, \tuple{3,0}, \dots
\]
దీనితో కింది ఫలితం స్థాపితమవుతుంది.
\end{explain}

\begin{prop}\ollabel{natsquaredenumerable}
$\Nat \times \Nat$ \tetoken{లెక్కించదగిన}{!!{enumerable}}.
\end{prop}

\begin{proof}
ప్రతి $k \in \Nat$ను, కాంటర్ జిగ్‌జాగ్ పట్టికలోని $n$వ అడ్డువరుస,
$m$వ నిలువువరుస స్థానంలో విలువ $k$ అయ్యే $\tuple{n,m} \in
\Nat \times \Nat$కు పంపేలా $f \colon \Nat \to \Nat\times\Nat$ను
తీసుకుందాం.
\end{proof}

\begin{explain}
ఈ పద్ధతి చక్కగా సాధారణీకృతమవుతుంది. ఉదాహరణకు, సహజ సంఖ్యల క్రమిత
త్రికాల సమితిని, అంటే
\[
\Nat \times \Nat \times \Nat = \Setabs{\tuple{n,m,k}}{n,m,k \in \Nat}
\]
ను లెక్కించడానికి దీన్ని వాడవచ్చు. $\Nat \times \Nat \times \Nat$ను
$\Nat \times \Nat$, $\Nat$ల కార్టీజియన్ లబ్ధంగా భావిస్తాం; అంటే,
\[
\Nat^3 = (\Nat \times \Nat) \times \Nat =
\Setabs{\tuple{\tuple{n,m},k}}{n, m, k
  \in \Nat }
\]
అందువల్ల, ఒక అక్షానికి $\Nat$ లెక్కింపునూ, మరొక అక్షానికి
$\Nat^2$ లెక్కింపునూ పేర్లుగా ఇచ్చి, ఒక పట్టికతో $\Nat^3$ను
లెక్కించవచ్చు:
\[
\begin{array}{ c | c | c | c | c | c}
& \mathbf 0 & \mathbf 1 & \mathbf 2 & \mathbf 3 & \dots \\
\hline
\mathbf{\tuple{0,0}} & \tuple{0,0,0} & \tuple{0,0,1} & \tuple{0,0,2} & \tuple{0,0,3} & \dots \\
\hline
\mathbf{\tuple{0,1}} & \tuple{0,1,0} & \tuple{0,1,1} & \tuple{0,1,2} & \tuple{0,1,3} & \dots \\
\hline
\mathbf{\tuple{1,0}} & \tuple{1,0,0} & \tuple{1,0,1} & \tuple{1,0,2} & \tuple{1,0,3} & \dots \\
\hline
\mathbf{\tuple{0,2}} & \tuple{0,2,0} & \tuple{0,2,1} & \tuple{0,2,2} & \tuple{0,2,3} & \dots\\
\hline
\vdots & \vdots & \vdots & \vdots & \vdots & \ddots \\
\end{array}
\]
ఇలా, కాంటర్ జిగ్‌జాగ్ పద్ధతి వంటి పద్ధతితో $\Nat^3$కు కూడా
లెక్కింపును పొందవచ్చు. ఇదే విధంగా కొనసాగిస్తూ, ఏ సహజ సంఖ్య $n$కైనా
$\Nat^n$ లెక్కింపును పొందవచ్చు. అందుచేత:
\end{explain}

\begin{prop}
ప్రతి $n \in \Nat$కు $\Nat^n$ \tetoken{లెక్కించదగిన}{!!{enumerable}}.
\end{prop}

\begin{prob}\label{sfr:siz:zigzag:prob:posint-n}
ప్రతి $n \in \Nat$కు $(\PosInt)^n$ \tetoken{లెక్కించదగిన}{!!{enumerable}} అని చూపండి.
\end{prob}

\begin{prob}\label{sfr:siz:zigzag:prob:posint-star}
  $(\PosInt)^*$ \tetoken{లెక్కించదగిన}{!!{enumerable}} అని చూపండి. మీరు
  \cref{sfr:siz:zigzag:prob:posint-n}ను వాడుకోవచ్చు.
\end{prob}

\end{document}
